\section{Discussion}\label{sec:disc}

\textbf{The competence is probably missing from the training data.}
VLA policies imitate demonstrations selected for task completion.
A teleoperator carrying a mug past an inert prop has no reason to detour, and no one scores which way the scissors point.
If the competence is absent from the training distribution, a safety instruction has nothing to retrieve and can trigger a related skill instead: the spill clause tilts the mug, plausibly toward a pour (\S\ref{sec:results}).
Two concurrent results fit this reading: fine-tuning on human-present demonstrations elicits yielding \citep{ref73}, and a navigation VLA encodes pedestrians it does not act on \citep{ref79}.
Our ablations are consistent with this hypothesis but do not establish it.

\textbf{The policy sets the demand on the safety layer.}
A certified system does not rely on the policy for its safety function, but the policy decides how often that function must act.
With a whole-arm protective stop, \pifive{} never reaches the crossing hand and still delivers the mug, but the stop fires on every carry and holds the arm for a median of 2.9\,s.
Some failures have no external fix at all: a stop does not change which way a blade points or how far a cup tilts.

\textbf{What a new policy should report.}
A low unsafe rate alone is not evidence of safety.
A policy that claims execution-phase safety should report each sub-type on matched cells against a person-blind control, next to its carry rate.
It should also show, with a feasibility witness, that the predicate can be met in the same scene.
Most of our comparisons could not have shown a policy safer, so larger matched cells come first.

\section{Limitations}\label{sec:limits}

{\setlength{\leftmargini}{1.1em}%
\begin{itemize}\setlength{\itemsep}{0pt}\setlength{\parskip}{0pt}\setlength{\topsep}{0pt}
\item \textbf{Simulation and proxies.} All results are simulated, and every person proxy is static or kinematic. There is no injury model, so contact and force counts are exposure, not harm.
\item \textbf{A capsule, not a human.} In the scored cells the policies see a capsule body. The human-mesh check covers T3 alone, on small samples, so the person-blind result may partly be blindness to a capsule.
\item \textbf{Small cells.} A tabletop cell holds eight episodes, so only large contrasts have power.
\item \textbf{One training corpus.} All four policies are trained on DROID, and rows differ in more than one of backbone, decoder and data, so no row difference is attributed to one of them.
\item \textbf{Reference-dependent T1, thin T2.} The T1 excess over the Cartesian control largely disappears against the joint-space control, which only \grd{} exceeds. T2 rests on one pre-registered placement with a thin margin.
\item \textbf{Unequal protocols.} \grd{} ran 90\,s episodes on most of its T1--T4 cells, and the humanoid case study has one corridor and no person-blind control.
\end{itemize}}
\noindent Finally, we show no learned policy that passes the benchmark; every feasibility witness is scripted or relies on an external stop.

\section{Conclusion}\label{sec:concl}

VLA safety evaluations ask whether a task should be done and whether it ended well; we measure how it is carried out, in four dimensions referenced to a nearby person.
Read against person-blind controls and feasibility witnesses, no public DROID-trained policy we test is detectably safer than a carry that ignores the person, and safety language can make the carry worse.
What a policy owns is not the safety function but the demand it places on that function, which can now be counted dimension by dimension.
Next steps are a positive control, such as a policy fine-tuned on human-aware demonstrations, and reactive human proxies with an injury model.
